\documentclass[11pt,a4paper]{article}

% ---------------------------------------------------------
% PREAMBLE (stable, warning-free, Overleaf-safe)
% ---------------------------------------------------------
\usepackage[utf8]{inputenc}
\usepackage[T1]{fontenc}
\usepackage{lmodern}
\usepackage{amsmath, amssymb, amsfonts}
\usepackage{bm}
\usepackage{geometry}
\usepackage{graphicx}
\usepackage{hyperref}
\usepackage{cite}
\usepackage{authblk}
\usepackage{physics}
\usepackage{mathtools}

\geometry{margin=2.4cm}

\title{\textbf{Companion LXXI: Hilbert-Space Embeddings of Persistence Diagrams in the Relaxation-Driven Cyclic Cosmology}}
\author{Michael Lehmann \\ \texttt{mi.lehmann@gmx.de}}
\date{April 2026}

\begin{document}
\maketitle

\begin{abstract}
Hilbert-space embeddings of persistence diagrams provide a linearized 
representation of topological structure, enabling inner-product computations, 
kernel methods, and spectral analyses. In weak-lensing analyses, such embeddings 
capture how connected components, loops, and void-like structures contribute to 
the geometry of topological feature space. In the standard cosmological model, 
these embeddings reflect the nonlinear matter distribution and the projection of 
large-scale structure. In the Relaxation-Driven Cyclic Cosmology (RDCC), the 
scale-dependent evolution of the gravitational potential modifies the Hilbert-space 
geometry of persistence diagrams in a characteristic way. This Companion develops 
a continuous description of Hilbert-space embeddings in RDCC, deriving how the 
relaxation scale influences linearized topology, kernel structure, and the 
observational signatures in weak-lensing surveys. The resulting framework provides 
a distinctive probe of the RDCC gravitational field through linearized topological 
geometry.
\end{abstract}
% ---------------------------------------------------------
\section{Introduction}

Persistence diagrams encode the birth and death of topological features in 
weak-lensing convergence fields. While wavelet transforms provide a multi-scale 
representation of topological fluctuations, Hilbert-space embeddings offer a 
linearized representation that enables inner-product computations, kernel methods, 
and spectral analyses.

In the standard cosmological model, the Hilbert-space geometry reflects the 
nonlinear matter distribution and the projection of large-scale structure. In 
RDCC, the gravitational potential evolves differently. The relaxation mechanism 
suppresses the decay of small-scale modes and enhances the coherence of large-scale 
modes. This modifies the Hilbert-space geometry of persistence diagrams in a 
scale-dependent manner, producing a distinctive signature in embedding analyses.
% ---------------------------------------------------------
\section{Hilbert-Space Embeddings of Persistence Diagrams}

A persistence diagram $D$ can be embedded into a reproducing kernel Hilbert space 
(RKHS) $\mathcal{H}$ via

\begin{equation}
    \Phi(D)
    = \sum_{(b,d) \in D}
      w(b,d) \, k(\cdot, (b,d)),
\end{equation}

where $k$ is a positive-definite kernel and $w(b,d)$ is a weight function.

The inner product becomes

\begin{equation}
    \langle \Phi(D_{1}), \Phi(D_{2}) \rangle_{\mathcal{H}}
    = \sum_{x \in D_{1}}
      \sum_{y \in D_{2}}
      w(x) w(y) k(x,y).
\end{equation}
% ---------------------------------------------------------
\section{RDCC Modification of Persistence Diagrams}

In RDCC, the convergence evolves as
\begin{equation}
    \kappa_{\mathrm{RDCC}}(\ell)
    = \kappa(\ell)
      \left[
        1 - \chi_{\kappa}
        \left(\frac{\ell}{\ell_{\mathrm{rel}}}\right)^{\alpha}
      \right].
\end{equation}

This modifies the birth and death thresholds of topological features:

\begin{align}
    b_{\mathrm{RDCC}} &= b \left[ 1 - \chi_{b} (\ell/\ell_{\mathrm{rel}})^{\alpha} \right], \\
    d_{\mathrm{RDCC}} &= d \left[ 1 - \chi_{d} (\ell/\ell_{\mathrm{rel}})^{\alpha} \right].
\end{align}

Thus the RDCC embedding becomes

\begin{equation}
    \Phi_{\mathrm{RDCC}}(D)
    = \Phi(D)
      \left[
        1 - \chi_{\Phi}
        \left(\frac{\ell}{\ell_{\mathrm{rel}}}\right)^{\alpha}
      \right].
\end{equation}
% ---------------------------------------------------------
\section{Kernel Structure in RDCC}

The RDCC-modified kernel is

\begin{equation}
    k_{\mathrm{RDCC}}(x,y)
    = k(x,y)
      \left[
        1 - \chi_{k}
        \left(\frac{\ell}{\ell_{\mathrm{rel}}}\right)^{\alpha}
      \right].
\end{equation}

This modification reflects the relaxation-driven change in linearized topology.
% ---------------------------------------------------------
\section{Linearized Topological Signatures of RDCC}

The relaxation mechanism enhances the coherence of large-scale modes. Hilbert-space 
embeddings therefore exhibit:

\begin{itemize}
    \item enhanced inner products for large-scale topological features,
    \item suppressed inner products for small-scale features,
    \item a characteristic turnover near $\ell_{\mathrm{rel}}$,
    \item modified kernel eigenvalue spectra,
    \item altered alignment of embedded persistence lifetimes.
\end{itemize}

These signatures provide a direct probe of the RDCC gravitational field through 
linearized topological geometry.
% ---------------------------------------------------------
\section{Conclusion}

Hilbert-space embeddings of persistence diagrams in the Relaxation-Driven Cyclic 
Cosmology reflect the scale-dependent evolution of the gravitational potential and 
the nonlinear matter distribution. The relaxation mechanism modifies the 
linearized topology in a scale-dependent manner, producing distinctive signatures 
in weak-lensing surveys. These effects provide a direct observational window into 
the relaxation scale and the temporal structure of the RDCC gravitational field. 
The present Companion establishes the theoretical foundation for future studies of 
kernel topology, nonlinear structure, and the interplay between light and gravity 
in RDCC.

\bibliographystyle{apsrev4-2}
\nocite{*}
\bibliography{references}


\end{document}
